\section{The numerical method, gently}
\label{sec:numerics}

The heat equation (\autoref{sec:heateq}) is a continuous statement about every
point and every instant. A computer can't handle ``infinitely many points.'' So
we \textbf{discretize}: chop space into finite slices and time into finite steps,
and turn the calculus into arithmetic the computer can repeat millions of times.

\subsection{Slicing the depth: why a \emph{geometric} grid}
\label{sec:grid}

We split the 5-metre soil column into many thin horizontal layers and track one
temperature per layer. The clever part: the layers are \textbf{not equally
thick}.
\begin{itemize}[leftmargin=1.4em,itemsep=3pt]
  \item Near the surface, temperature swings violently over a few centimetres
    (the skin). To capture that we need \textbf{very thin} layers
    (${\sim}\SI{2}{mm}$ at the top).
  \item Deep down, temperature barely changes over a whole metre. Thin layers
    there would just waste computer effort.
\end{itemize}
So each layer is made ${\sim}8\%$ thicker than the one above it. This is a
\textbf{geometric grid} (each step multiplied by a constant ratio).

\begin{figure}[h]
\centering
\begin{tikzpicture}[x=1cm,y=1cm]
  % geometric layer boundaries (exaggerated for clarity)
  \def\xL{0} \def\xR{4}
  \pgfmathsetmacro{\prev}{0}
  \foreach \i/\h in {1/0.18, 2/0.26, 3/0.36, 4/0.50, 5/0.70, 6/0.98, 7/1.36, 8/1.86} {
    \pgfmathsetmacro{\top}{\prev}
    \pgfmathsetmacro{\bot}{\prev+\h}
    \fill[teal!\the\numexpr8+\i*8\relax] (\xL,-\top) rectangle (\xR,-\bot);
    \draw[char,line width=0.4pt] (\xL,-\top) rectangle (\xR,-\bot);
    \node[char,font=\scriptsize] at (\xR+1.5,{-(\top+\bot)/2}) {layer \i};
    \global\let\prev\bot
  }
  % surface label + sun
  \node[coral,font=\bfseries\small,anchor=south] at (2,0.15) {surface (sunlight in)};
  \draw[-{Latex[length=2mm]},coral,line width=1pt] (-0.9,0.9) -- (-0.3,0.15);
  % depth arrow
  \draw[-{Latex[length=2.5mm]},char,line width=0.9pt] (-0.7,0) -- (-0.7,-6.0)
    node[midway,left,rotate=90,anchor=south,font=\small] {increasing depth $z$};
  % bracket on first thin layer
  \draw[decorate,decoration={brace,amplitude=4pt},dim]
    (\xR+0.15,0) -- (\xR+0.15,-0.18) node[midway,right,font=\scriptsize,xshift=2pt] {${\sim}2$ mm};
\end{tikzpicture}
\caption[The geometric depth grid]{The geometric depth grid (thicknesses
exaggerated). \textbf{What to notice:} slices are paper-thin at the top, where
the daily heat wave is sharp, and grow ${\sim}8\%$ thicker with each step
downward, so the deep, sleepy soil is covered cheaply. \textbf{Why it matters:}
this puts the computational effort exactly where the physics is fast-changing
--- like taking many photos per second of a sprint start but only one per minute
of a marathon's middle.}
\label{fig:grid}
\end{figure}

The grid docstring in \file{lunar/grid.py} puts it perfectly:

\begin{lstlisting}[language={},style=shellstyle,caption={From the \file{lunar/grid.py} docstring.}]
Near the surface, where temperature swings wildly between lunar day
and night, we use paper-thin slices (~2 mm) so we capture the fast
changes. Deeper down, where temperature barely moves, the slices get
gradually thicker (each ~8% thicker than the one above). This
"geometric" spacing puts the computational effort where the action is.
\end{lstlisting}

The grid is built by \code{make\_geometric\_grid}, configured in
\file{lunar/config.py}:

\begin{lstlisting}[caption={The production grid in \file{lunar/config.py}.}]
# --- depth grid (geometric) ---
GRID = dict(z_max=5.0, dz0=0.002, growth=0.08)
\end{lstlisting}

That reads: total depth \SI{5}{m}, top layer \SI{2}{mm} thick, each next layer
$8\%$ thicker (growth ratio 1.08). Uniform grids are \emph{deliberately
forbidden} in this project (the code even raises an error) because they
under-resolve the skin while wasting effort deep down.

\begin{headsup}
\file{lunar/constants.py} also defines \emph{fallback} grid defaults
(\code{Z\_MAX\_DEFAULT = 3.0}, \code{GROWTH\_DEFAULT = 0.15}). Those are only used
if someone calls \code{make\_geometric\_grid()} with no arguments. The real
science runs always pass \code{GRID} from \file{config.py} (\SI{5}{m} /
\SI{2}{mm} / 1.08), so the production grid is the \file{config.py} one. Don't be
confused by the two different numbers.
\end{headsup}

\subsection{Stepping through time: what a ``timestep'' is}
\label{sec:timestep}

We also chop \textbf{time} into steps. Here each step $\Delta t$ is \textbf{one
hour} (\code{DT\_STEP = 3600.0} seconds in \file{config.py}). Starting from the
temperatures right now, we apply the heat equation to compute the temperatures
one hour later, then repeat. One full lunar day (a ``lunation'') is
${\sim}29.5$ Earth days (\SI{2551443}{s}), so one simulated lunation is
\textbf{709 hourly samples} (\code{N\_t = int(T\_LUNAR / DT\_STEP) + 1} in
\code{run\_with}). The test suite samples more coarsely (2-hour steps) to keep
continuous-integration runs fast.

\begin{headsup}
\textbf{A handy units gotcha.} Watch out for the number 2551: it looks like a
timestep count but is really the lunation \emph{length} in thousands of seconds
(\SI{2551443}{s}). The number of hourly steps is \textbf{709}.
\end{headsup}

\subsection{Crank--Nicolson: a stable, accurate way to take a step}
\label{sec:cn}

There are simple ways to step a heat equation forward, but the na\"ive ones
either blow up (become numerically unstable) unless the timestep is
microscopically small, or they are inaccurate. \textbf{Crank--Nicolson} is the
standard professional choice. Intuition:
\begin{itemize}[leftmargin=1.4em,itemsep=3pt]
  \item An \textbf{explicit} step computes the future purely from the present.
    Cheap, but unstable for stiff problems like heat (it can oscillate and
    explode).
  \item An \textbf{implicit} step computes the future from the future (you solve
    a small system of equations). Stable, but slightly less accurate.
  \item \textbf{Crank--Nicolson averages the two} (``half now, half later,'' the
    $0.5$ factors you'll see in the code). The result is both \textbf{stable}
    \emph{and} \textbf{second-order accurate} in time. It is coded in
    \code{\_step} in \file{lunar/solver.py} (look for the \code{0.5 * alpha\_...}
    terms --- that $0.5$ \emph{is} the Crank--Nicolson averaging).
\end{itemize}

\begin{keyidea}
You don't need to derive it. Just know: Crank--Nicolson is the recipe that lets
us take comfortable one-hour steps without the simulation exploding.
\end{keyidea}

\subsection{The tridiagonal (Thomas) solver: the fast inner engine}
\label{sec:thomas}

An implicit / Crank--Nicolson step requires solving a system of linear
equations: ``find all the new-layer temperatures at once such that they are
mutually consistent.'' Normally solving $N$ equations in $N$ unknowns is slow
(cost grows like $N^3$). But here each layer only talks to its \textbf{immediate
neighbours above and below} --- so the system is \textbf{tridiagonal} (nonzero
only on the diagonal and the two adjacent diagonals).

For tridiagonal systems there is a famous shortcut, the \textbf{Thomas
algorithm}, that solves them in a single sweep down and a single sweep back up
--- cost grows only like $N$ (linear), which is as fast as it gets. It is
implemented in \code{\_thomas} in \file{lunar/solver.py}:

\begin{lstlisting}[caption={The Thomas tridiagonal solver, \file{lunar/solver.py}.}]
@njit(cache=True)
def _thomas(a, b, c, d):
    """Thomas algorithm for an (a, b, c) tridiagonal system A x = d."""
    n = b.size
    cp = np.empty(n); dp = np.empty(n); x = np.empty(n)
    cp[0] = c[0] / b[0]
    dp[0] = d[0] / b[0]
    for i in range(1, n):                 # forward sweep
        m = b[i] - a[i] * cp[i - 1]
        cp[i] = c[i] / m if i < n - 1 else 0.0
        dp[i] = (d[i] - a[i] * dp[i - 1]) / m
    x[n - 1] = dp[n - 1]
    for i in range(n - 2, -1, -1):        # back substitution
        x[i] = dp[i] - cp[i] * x[i + 1]
    return x
\end{lstlisting}

The \code{@njit} decorator is \textbf{Numba}, a tool that compiles this Python
down to machine code so it runs ${\sim}10$--$100\times$ faster (this inner loop
runs millions of times). If Numba isn't installed, the code falls back to plain
Python and still works --- see the \code{try/except} at the top of
\file{lunar/solver.py}. (\file{docs/} contains \file{NUMBA\_VS\_CPP\_JUSTIFICATION.md}
explaining why Numba was chosen over writing C++.)

\begin{figure}[h]
\centering
\includegraphics[width=0.82\textwidth]{fig_phase4_njit_vs_cpp.png}
\caption[Numba speed-up]{Why the \code{@njit} compiler matters: a benchmark of
the solver's inner loop. \textbf{What to notice:} compiling the hot Python loop
with Numba closes almost all of the gap to hand-written C++, while keeping the
code in one readable language. \textbf{Why it matters:} the retrieval runs the
solver hundreds of times, so this speed-up is what makes the whole study
practical on a laptop.}
\label{fig:njit}
\end{figure}

One more detail: conductivity between two adjacent layers uses the
\textbf{harmonic mean} (\code{\_face\_harmonic\_mean}), not the simple average.
For heat flowing \emph{through} two materials in series (like the deep,
conductive layer feeding the fluffy, insulating layer above), the harmonic mean
is the physically correct combination --- the poor conductor dominates, just as
the narrowest pipe dominates a series of pipes.

\subsection{Newton's method for the non-linear surface}
\label{sec:newton}

Remember the surface radiates as $T_s^4$ (\autoref{sec:boundaries}). That fourth
power makes the surface energy-balance equation \textbf{non-linear} --- you can't
just solve for $T_s$ algebraically. Instead we use \textbf{Newton's method}, the
classic ``guess, check the error, slide toward the answer, repeat'' root-finder:
\begin{enumerate}[leftmargin=1.5em,itemsep=2pt]
  \item Guess a surface temperature $T_s$.
  \item Compute the \textbf{residual} $R(T_s)$ --- how badly the energy balance
    fails (absorbed $-$ radiated $-$ conducted;
    \code{surface\_energy\_balance\_residual}).
  \item Use the slope $\mathrm{d}R/\mathrm{d}T_s$ to jump to a better guess.
  \item Repeat until $R$ is ${\sim}0$ (a few iterations).
\end{enumerate}
This is \code{\_solve\_surface\_newton} in \file{lunar/solver.py}. The code notes
the derivative is always negative, which guarantees Newton converges from any
positive starting temperature. So at \textbf{every hourly step}, the solver does
a tiny Newton solve to pin the surface temperature, then a Thomas sweep to update
the whole column. That's the complete inner loop.
